\documentclass[11pt,a4paper]{article}
\usepackage[T1]{fontenc}
\usepackage{lmodern,amsmath,amssymb,amsthm,booktabs,microtype}
\usepackage[margin=26mm]{geometry}
\usepackage[hidelinks,pdftitle={Which central elements are transitive?},pdfauthor={Anonymous}]{hyperref}
\usepackage{enumitem}
\setlength{\emergencystretch}{2em}
\newtheorem{theorem}{Theorem}[section]
\newtheorem{lemma}[theorem]{Lemma}
\newtheorem{proposition}[theorem]{Proposition}
\newtheorem{corollary}[theorem]{Corollary}
\theoremstyle{definition}
\newtheorem{definition}[theorem]{Definition}
\newtheorem{remark}[theorem]{Remark}
\newtheorem{question}[theorem]{Question}
\newcommand{\Sn}{\mathfrak S}
\newcommand{\Z}{\mathbb Z}
\newcommand{\Q}{\mathbb Q}
\newcommand{\C}{\mathbb C}
\newcommand{\one}{\mathbf 1}
\newcommand{\Pa}{\Pi}
\DeclareMathOperator{\cyc}{cyc}
\DeclareMathOperator{\Orb}{Orb}
\DeclareMathOperator{\type}{type}
\title{Which central elements are transitive?\\[4pt]
\large An answer to a question of Campion Loth and Rattan on Jucys--Murphy evaluations}
\author{Anonymous}
\date{November 2026\\[2pt]\small DOI \href{https://doi.org/10.5281/zenodo.24516957}{10.5281/zenodo.24516957}}
\begin{document}
\maketitle

\begin{abstract}
Symmetric polynomials in the Jucys--Murphy elements of the symmetric group $\Sn_n$ give exactly the centre of its group algebra. Campion Loth and Rattan showed that the elements stay central if, after expanding such a polynomial into products of transpositions, one keeps only the products whose transpositions connect all $n$ letters. They asked which subalgebra of the centre these transitive images generate (their Question~12). We show that over the complex numbers no generation is needed: every central element is the transitive image of a single symmetric function. Over an arbitrary commutative coefficient ring, the transitive images span the centre, and also generate it as a unital algebra, exactly when $(n-1)!$ is invertible. When they span it, every central element has a preimage of degree at most $2n-2$, and this degree cannot be lowered. We also describe which pairs consisting of an ordinary and a transitive evaluation of the same function occur. Over the integers, a conjugation-orbit argument bounds the image classwise. An exact, degree-free computation shows that this bound describes the integral image exactly for $n\le 5$ and $n=7$, but not for $n=6$, and that adding the scalars always closes the image under multiplication for $n\le7$.
\end{abstract}

\section{Introduction}\label{sec:intro}
This section explains the question, why it is natural, what was known before, and what this note proves.

\paragraph{The objects.} The \emph{Jucys--Murphy elements} of the symmetric group $\Sn_n$ are the sums of transpositions
\[
 J_1=0,\qquad J_i=(1\ i)+(2\ i)+\cdots+(i-1\ i)\quad(2\le i\le n)
\]
in the group algebra. They commute with one another, and a classical theorem of Jucys and of Murphy states that the symmetric polynomials in $J_1,\ldots,J_n$ are exactly the central elements of the group algebra, even over the integers \cite{jucys,murphy}; see also \cite[Theorem~2]{clr} and \cite[Introduction]{murray}. Each $J_i$ is a sum of transpositions, so expanding a symmetric polynomial $f(J_1,\ldots,J_n)$ produces a sum of products of transpositions. The coefficient of a permutation $\sigma$ in such a sum counts factorisations of $\sigma$ into transpositions of a prescribed kind. Monotone Hurwitz numbers and star factorisation numbers arise in this way \cite{clr,gj}.

\paragraph{Transitivity.} In enumerative geometry one usually wants \emph{transitive} factorisations, those whose transpositions connect all $n$ letters; they correspond to connected branched covers of the sphere. Keeping only the transitive products in the expansion of $f(J_1,\ldots,J_n)$ defines the \emph{transitive image} $T_n(f)$ (Definition~\ref{def:T}). Goulden and Jackson proved that the transitive powers $T_n(p_k)$ are central, which explains a surprising symmetry of star factorisations \cite{gj}. F\'eray gave a conceptual proof of that centrality through partial permutations \cite{feray}, and Campion Loth and Rattan proved that $T_n(f)$ is central for every symmetric function $f$ \cite[Theorem~6]{clr}. They then asked the following.

\begin{question}[{\cite[Question~12]{clr}}]
What is the subalgebra of $Z(\C\Sn_n)$ generated by the transitive images $T_n(f)$ of all symmetric functions $f$?
\end{question}

The question is natural because $T_n$ is linear but not multiplicative: a product of two transitive words is transitive, but a transitive word can be a product of two intransitive ones. It is therefore not clear whether discarding the disconnected factorisations loses information about the centre.

\paragraph{Results.} We answer the question, and its analogue over any commutative ring $R$ with identity, as follows.
\begin{enumerate}[label=(\Alph*),leftmargin=2.2em]
\item \emph{Linear image} (Theorem~\ref{th:centre}). The map $T_n$ from symmetric functions over $R$ to $Z(R\Sn_n)$ is onto if and only if $(n-1)!$ is a unit in $R$. In particular, over $\C$ every central element is a single transitive image, so the subalgebra asked for in Question~12 is the whole centre. Over every $R$, the coefficients of the identity permutation in transitive images form exactly the ideal $(n-1)!\,R$.
\item \emph{Optimal degree} (Theorem~\ref{th:degree}). When $(n-1)!$ is a unit and $R\ne0$, every central element is $T_n(f)$ for some $f$ of degree at most $2n-2$, and the identity element is not $T_n(f)$ for any $f$ of smaller degree.
\item \emph{Ordinary and transitive evaluations together} (Theorem~\ref{th:pair}). When $(n-1)!$ is a unit, the pairs $(f(J),T_n(f))$ are exactly the pairs of central elements with equal coefficients at the $n$-cycles.
\item \emph{Algebra generation and integrality} (Theorem~\ref{th:unital}, Proposition~\ref{prop:orbit}, Theorem~\ref{th:integral}). Even the unital algebra generated by the transitive images is the whole centre only when $(n-1)!$ is a unit; for each prime $p\le n-1$ the class of a $p$-cycle stays out of reach modulo $p$. More generally, every element of that algebra has coefficient divisible by $z_\lambda/h_\lambda$ at every class $\lambda\ne(1^n)$, and every transitive image also at $\lambda=(1^n)$ (notation in Section~\ref{sec:setup}). An exact computation of the integral image for $n\le7$ shows that this bound is attained classwise for $n\le5$ and $n=7$, but that for $n=6$ the integral image has index $2$ in the lattice it allows.
\end{enumerate}
Part (A) answers Question~12 as posed. Parts (B)--(D) describe the answer over other coefficient rings and the cost of using transitive images, in degree and in arithmetic.

\paragraph{Method.} Everything rests on a M\"obius inversion over set partitions (Lemma~\ref{lem:mob}), which writes $T_n(f)$ as an alternating sum of ordinary evaluations of $f$ on Jucys--Murphy elements of Young subgroups. Inclusion--exclusion over subsets was already used by Goulden and Jackson for star powers \cite[Section~1.1]{gj}, and passing between all and connected objects by M\"obius inversion is standard in Hurwitz theory. The new ingredient is the choice of symmetric functions on which this formula is easy to evaluate: multiplying by the sum of the elements of a Young subgroup collapses every Jucys--Murphy element to the number of smaller letters in its block. This lets one read off the block sizes of a set partition and invert the resulting triangular system.

\paragraph{What is not done here.} We do not find closed forms for transitive images of the classical bases (Question~13 of \cite{clr}), and we do not determine the integral image for every $n$; Section~\ref{sec:integral} states what remains open.

\paragraph{Roadmap.} Section~\ref{sec:setup} fixes notation and proves the M\"obius formula. Section~\ref{sec:linear} proves the surjectivity criterion. Section~\ref{sec:degree} proves the optimal degree, which improves the degree bound $n^2$ of the construction in Section~\ref{sec:linear} to $2n-2$. Section~\ref{sec:pair} treats ordinary and transitive evaluations together. Section~\ref{sec:integral} treats algebra generation and the integral image. Section~\ref{sec:checks} describes the computer checks, and Section~\ref{sec:remarks} adds remarks and the AI-use statement.

\section{Notation and the M\"obius formula}\label{sec:setup}
This section fixes notation, defines the transitive image carefully, and proves the formula on which all later arguments rest.

\paragraph{Groups and classes.} Fix $n\ge1$, write $[n]=\{1,\ldots,n\}$, and let $R$ be a commutative ring with identity. Products of permutations may be read in either order; no statement below depends on the convention, because under the other convention a word has the inverse product and the same orbit partition, and a permutation and its inverse have the same cycle type. For $x\in R\Sn_n$ write $[\sigma]x$ for the coefficient of $\sigma$, and $\varepsilon(x)=\sum_\sigma[\sigma]x$ for the \emph{augmentation}, a ring homomorphism $R\Sn_n\to R$. Write $\cyc(\sigma)$ for the number of cycles of $\sigma$, fixed points included. For an integer partition $\lambda\vdash n$ write $\ell(\lambda)$ for its number of parts, $m_i(\lambda)$ for the number of parts equal to $i$,
\[
 z_\lambda=\prod_i i^{m_i(\lambda)}m_i(\lambda)!,\qquad h_\lambda=\gcd\{\,i\,m_i(\lambda): m_i(\lambda)>0\,\},
\]
and $K_\lambda$ for the sum of the permutations of cycle type $\lambda$. The class of type $\lambda$ has $n!/z_\lambda$ elements. The class sums form a basis of the centre $Z(\Z\Sn_n)$ over $\Z$, hence of $Z(R\Sn_n)$ over $R$.

\paragraph{Set partitions.} Let $\Pa_n$ be the set of set partitions of $[n]$, ordered by refinement: $\xi\le\pi$ when every block of $\xi$ lies inside a block of $\pi$. Its largest element is $\hat1=\{[n]\}$. The \emph{type} of $\pi$ is the integer partition formed by its block sizes. Put $\mu(\pi)=(-1)^{k-1}(k-1)!$ when $\pi$ has $k$ blocks.

\begin{lemma}\label{lem:mu}
For every $\xi\in\Pa_n$, $\sum_{\pi\ge\xi}\mu(\pi)=1$ if $\xi=\hat1$ and $0$ otherwise.
\end{lemma}
\begin{proof}
Partitions $\pi\ge\xi$ correspond to set partitions of the $m$ blocks of $\xi$, with the same number of blocks, so the sum equals $s_m=\sum_kS(m,k)(-1)^{k-1}(k-1)!$, where $S(m,k)$ counts set partitions of an $m$-set into $k$ blocks. By the exponential formula, $\sum_{m\ge1}s_mx^m/m!=\sum_{k\ge1}(-1)^{k-1}(e^x-1)^k/k=\log(1+(e^x-1))=x$. Hence $s_1=1$ and $s_m=0$ for $m\ge2$.
\end{proof}
Thus $\mu(\pi)$ is the M\"obius function $\mu(\pi,\hat1)$ of the partition lattice.

\paragraph{Symmetric functions.} Let $\Lambda_R$ be the ring of symmetric functions over $R$; the products $e_\nu=e_{\nu_1}e_{\nu_2}\cdots$ of elementary symmetric functions form an $R$-basis, and $\Lambda_R=R\otimes_\Z\Lambda_\Z$. For $f\in\Lambda_R$ write $E_n(f)=f(J_1,\ldots,J_n,0,0,\ldots)$, the \emph{ordinary evaluation}. By the theorem of Jucys and Murphy, $E_n(\Lambda_\Z)=Z(\Z\Sn_n)$: Murphy proves that the centre of $K\Sn_n$ is the ring of symmetric polynomials over $K$ in $J_2,\ldots,J_n$ for every integral domain $K$ \cite[Theorem~1.9]{murphy}, and Jucys obtains it from his identity \eqref{eq:ek} for $\pi=\hat1$ and the Farahat--Higman generators of the integral centre, noting that the result holds over the integers \cite[Sections~2--3]{jucys}; by base change $E_n(\Lambda_R)=Z(R\Sn_n)$ for every $R$.

\paragraph{Words and the transitive image.} A \emph{word} is a finite sequence $w=(t_1,\ldots,t_k)$ of transpositions; its product is $\overline w=t_1\cdots t_k$, and its \emph{orbit partition} $\Orb(w)\in\Pa_n$ is the partition into connected components of the graph on $[n]$ whose edges are the $t_r$. The word is \emph{transitive} when $\Orb(w)=\hat1$. An \emph{ordered monomial} is a sequence $(i_1,\ldots,i_k)$ of indices; substituting $J_{i_r}=\sum_{a<i_r}(a\ i_r)$ expands it into the $\prod_r(i_r-1)$ words $((a_1\ i_1),\ldots,(a_k\ i_k))$ with $a_r<i_r$.

\begin{definition}\label{def:T}
Write $f(x_1,\ldots,x_n,0,\ldots)$ as an $R$-linear combination $\sum_mc_m\,x_{m}$ of ordered monomials $x_m=x_{i_1}\cdots x_{i_k}$ (any such expression; repetitions allowed). The \emph{transitive image} is
\[
 T_n(f)=\sum_mc_m\sum_{w}\overline w\in R\Sn_n,
\]
where the inner sum runs over the transitive words in the expansion of the ordered monomial $m$.
\end{definition}

This is the operator of Campion Loth and Rattan \cite[Section~1.3]{clr}. Because $T_n$ is not a map on the group algebra, its independence of the chosen expression has to be proved; Lemma~\ref{lem:mob} does this. For $\pi\in\Pa_n$ put
\[
 J_i^\pi=\sum_{\substack{a<i\\ a,i\text{ in the same block of }\pi}}(a\ i),\qquad
 \Sn_\pi=\prod_{B\in\pi}\Sn_B,\qquad H_\pi=\sum_{\sigma\in\Sn_\pi}\sigma .
\]
For $i$ in a block $B$, $J^\pi_i$ is the Jucys--Murphy element of $\Sn_B$ attached to $i$ after relabelling $B$ in increasing order, and transpositions in different blocks commute. Hence $J_1^\pi,\ldots,J_n^\pi$ commute, $f(J^\pi)$ is well defined, and $J^\pi_i=0$ when $i$ is the least letter of its block; in particular $J^\pi_1=0$.

\begin{lemma}[M\"obius formula]\label{lem:mob}
For every expression in Definition~\ref{def:T},
\begin{equation}\label{eq:mob}
 T_n(f)=\sum_{\pi\in\Pa_n}\mu(\pi)\,f(J^\pi).
\end{equation}
In particular $T_n(f)$ does not depend on the expression chosen, and $T_n:\Lambda_R\to R\Sn_n$ is $R$-linear.
\end{lemma}
\begin{proof}
Expand $f(J^\pi)$ using the same expression for $f$. The words that occur are exactly the words $w$ of the expansion of $f(J)$ all of whose transpositions lie inside blocks of $\pi$, that is, those with $\Orb(w)\le\pi$, each with the same coefficient $c_m$. Hence the right side of \eqref{eq:mob} equals $\sum_mc_m\sum_w\overline w\sum_{\pi\ge\Orb(w)}\mu(\pi)$, which is $T_n(f)$ by Lemma~\ref{lem:mu}. The right side does not depend on the expression, because each $f(J^\pi)$ is the evaluation of a polynomial at commuting elements.
\end{proof}

\paragraph{Block identities.} For a block $B$ of size $b$, write $J^B$ for the family $(J^\pi_i)_{i\in B}$. The following classical identity of Jucys is equivalent to \cite[equation~(8)]{jucys} and \cite[Theorem~1]{clr}:
\begin{equation}\label{eq:jucys}
 \prod_{i\in B}(u+J^\pi_i)=\sum_{\sigma\in\Sn_B}u^{\cyc_B(\sigma)}\sigma ,
\end{equation}
where $\cyc_B(\sigma)$ counts the cycles of $\sigma$ on $B$. (Proof: multiply the factors in increasing order of $i$. Multiplying a permutation of the smaller letters by $u$ makes $i$ a new fixed point, and multiplying by $(a\ i)$ inserts $i$ into the cycle of $a$; every permutation of the larger set arises once.) Comparing coefficients of powers of $u$ shows that the expansion of $e_k(J^B)$ contains each permutation $\omega\in\Sn_B$ with $b-\cyc_B(\omega)=k$ exactly once, as a word with strictly increasing larger letters. In that word each transposition joins two cycles of the partial product, since the number of cycles falls from $b$ to $b-k$ in $k$ steps of $\pm1$; by induction on the length, the orbit partition of the word equals the cycle partition of $\omega$. Multiplying over the blocks, $\prod_i(1+J_i^\pi)=H_\pi$, and
\begin{equation}\label{eq:ek}
 e_k(J^\pi)=\sum_{\substack{\omega\in\Sn_\pi\\ n-\cyc(\omega)=k}}\omega .
\end{equation}
Conjugation by $\tau\in\Sn_n$ maps $\Sn_\pi$ onto $\Sn_{\tau\pi}$ and preserves cycle numbers, so it carries $e_\nu(J^\pi)$ to $e_\nu(J^{\tau\pi})$. Since $\mu(\tau\pi)=\mu(\pi)$, formula \eqref{eq:mob} shows again that $T_n(f)$ is central for every $f\in\Lambda_R$, the theorem of Campion Loth and Rattan \cite[Theorem~6]{clr}, here over any $R$.

\paragraph{Transitive images as tuple counts.} Applying Definition~\ref{def:T} to $e_\nu=e_{\nu_1}\cdots e_{\nu_l}$, expanded factor by factor with strictly increasing larger letters inside each factor, the description above gives
\begin{equation}\label{eq:count}
 [\sigma]\,T_n(e_\nu)=\#\bigl\{(\omega_1,\ldots,\omega_l)\in\Sn_n^l:\ n-\cyc(\omega_r)=\nu_r,\ \omega_1\cdots\omega_l=\sigma,\ \langle\omega_1,\ldots,\omega_l\rangle\text{ transitive}\bigr\},
\end{equation}
as in the proof of \cite[Theorem~6]{clr}. In particular $T_n(\Lambda_\Z)\subseteq Z(\Z\Sn_n)$, and because $\Lambda_R=R\otimes\Lambda_\Z$, the image $T_n(\Lambda_R)$ is the $R$-span of $T_n(\Lambda_\Z)$. Note also that $e_k(J^\pi)=0$ for $k\ge n$, since $\Sn_\pi$ has no permutation with fewer than one cycle.

\section{The linear image}\label{sec:linear}
This section proves that transitive images fill the centre exactly when $(n-1)!$ is invertible. The sufficiency proof is constructive: it produces, for each type $\alpha$ of set partition, a symmetric function whose transitive image is a unit multiple of the sum of the Young-subgroup sums $H_\pi$ of that type.

For $\alpha\vdash n$ put $G_\alpha=\sum_{\type(\pi)=\alpha}H_\pi$.

\begin{lemma}\label{lem:G}
The elements $G_\alpha$, $\alpha\vdash n$, form a $\Z$-basis of $Z(\Z\Sn_n)$, related to the class sums by an integer unitriangular matrix.
\end{lemma}
\begin{proof}
Each $G_\alpha$ is invariant under conjugation, hence central. For $\sigma$ of cycle type $\beta$, the coefficient $[\sigma]G_\alpha$ is the number of set partitions $\pi$ of type $\alpha$ with $\sigma\in\Sn_\pi$, that is, the number of ways to group the cycles of $\sigma$ into blocks with sizes $\alpha$. It vanishes unless $\beta$ refines $\alpha$. If $\beta=\alpha$ it equals $1$: the grouping has as many blocks as $\sigma$ has cycles, so each block is a single cycle. Ordering partitions by any linear extension of refinement makes the matrix unitriangular.
\end{proof}

\begin{theorem}[Linear image]\label{th:centre}
\begin{enumerate}[label=(\roman*)]
\item If $(n-1)!$ is a unit in $R$, then $T_n:\Lambda_R\to Z(R\Sn_n)$ is onto. Explicitly, there are symmetric functions $F_\alpha$ of degree at most $n^2$ with $T_n(F_\alpha)=\mu_\alpha G_\alpha$, where $\mu_\alpha=(-1)^{\ell(\alpha)-1}(\ell(\alpha)-1)!$.
\item For every $R$, $\{[1]\,T_n(f): f\in\Lambda_R\}=(n-1)!\,R$.
\end{enumerate}
Consequently $T_n$ is onto $Z(R\Sn_n)$ if and only if $(n-1)!$ is a unit in $R$.
\end{theorem}
\begin{proof}
(i) For $0\le j<n$ let $L_j$ be the Lagrange polynomial with $L_j(k)=\delta_{jk}$ for $k=0,\ldots,n-1$; its denominator $j!\,(n-1-j)!$ is a unit. With independent indeterminates $y_0,\ldots,y_{n-1}$ (markers) put
\[
 F(x;y)=\prod_{i=1}^n\Bigl[(1+x_i)\sum_{j=0}^{n-1}y_jL_j(x_i)\Bigr],
\]
a symmetric polynomial in $x_1,\ldots,x_n$ with coefficients in $R[y]$. Fix $\pi$. For $i$ in a block $B$ let $r_i$ be the number of letters of $B$ smaller than $i$; then $\varepsilon(J^\pi_i)=r_i$, and the values $r_i$ for $i\in B$ are $0,1,\ldots,|B|-1$. Since $\prod_i(1+J^\pi_i)=H_\pi$ and $H_\pi v=\varepsilon(v)H_\pi$ for every $v\in R\Sn_\pi$, we get
\[
 F(J^\pi;y)=H_\pi\prod_i\sum_jy_jL_j(r_i)=H_\pi\prod_iy_{r_i}=H_\pi\prod_{j=0}^{n-1}y_j^{a_j(\pi)},\qquad a_j(\pi)=\#\{B\in\pi:|B|>j\}.
\]
The exponents determine the type of $\pi$, since $a_j-a_{j+1}$ is the number of blocks of size $j+1$. Let $F_\alpha$ be the coefficient in $F$ of the marker monomial $\prod_jy_j^{a_j}$ belonging to type $\alpha$, lifted to $\Lambda_R$ through its expression in elementary symmetric functions (which does not raise the degree, at most $n\cdot n$). By \eqref{eq:mob}, $T_n(F_\alpha)=\sum_{\type(\pi)=\alpha}\mu(\pi)H_\pi=\mu_\alpha G_\alpha$. The integer $\mu_\alpha$ divides $(n-1)!$, so it is a unit, and Lemma~\ref{lem:G} expresses every class sum as an integer combination of the $G_\alpha$.

(ii) By the end of Section~\ref{sec:setup} it suffices to treat $f=e_\nu$ over $\Z$. Let $X$ be the set of tuples in \eqref{eq:count} whose product is the identity. The group $\Sn_n$ acts on $X$ by simultaneous conjugation. A permutation $c$ fixing a tuple commutes with the transitive group $G$ it generates; if $c$ fixes a letter $x$ it fixes every $g(x)$, $g\in G$, so $c=1$. Hence the stabiliser acts freely on $[n]$, its order divides $n$, and every orbit has size divisible by $n!/n=(n-1)!$. So $[1]T_n(e_\nu)=|X|\in(n-1)!\,\Z$. Conversely, $e_{n-1}(J)=K_{(n)}$ by \eqref{eq:ek}, and the identity coefficient of $T_n(e_{n-1}e_{n-1})$ counts the pairs $(\omega,\omega^{-1})$ with $\omega$ an $n$-cycle, of which there are $(n-1)!$. (For $n=1$ the empty word is transitive and $T_1(1)=1$.) If $T_n$ is onto, then $1\in(n-1)!R$, so $(n-1)!$ is a unit.
\end{proof}

For example, for $n=3$, with columns ordered $K_{111},K_{21},K_3$, the three elements $\mu_\alpha G_\alpha$ for $\alpha=(1,1,1),(2,1),(3)$ are the rows of
\[
 \begin{pmatrix}2&0&0\\-3&-1&0\\1&1&1\end{pmatrix}.
\]
Here $F_{(1,1,1)}$ has denominator $2$; consistently with part (ii), no integral symmetric function has a transitive image with identity coefficient $1$ when $n=3$.

\section{The optimal degree}\label{sec:degree}
This section shows that degree $2n-2$ suffices for every central element and that no smaller degree reaches the identity. The construction replaces the Lagrange markers of Section~\ref{sec:linear} by polynomials of degree at most $n-1$ that recognise the type of a set partition.

\begin{theorem}[Optimal degree]\label{th:degree}
Let $R\ne0$ and let $(n-1)!$ be a unit in $R$. Every element of $Z(R\Sn_n)$ is $T_n(f)$ for some $f\in\Lambda_R$ of degree at most $2n-2$. If $f$ has degree less than $2n-2$, then $[1]\,T_n(f)=0$; in particular the identity element is not the transitive image of any symmetric function of degree less than $2n-2$.
\end{theorem}
\begin{proof}
\emph{Evaluation on blocks.} For a symmetric polynomial $g$, the computation in the proof of Theorem~\ref{th:centre} gives $H_\pi\,g(J^\pi)=g(X_\pi)H_\pi$, where $X_\pi$ is the multiset $\bigcup_{B\in\pi}\{0,1,\ldots,|B|-1\}$. For $1\le s\le n$ put $B_s(x)=\sum_{i=1}^n\binom{x_i}{s-1}$, a symmetric polynomial of degree $s-1$ whose denominator $(s-1)!$ is a unit. Summing $\binom{r}{s-1}$ over $r<|B|$ gives
\[
 B_s(X_\pi)=\sum_{B\in\pi}\binom{|B|}{s}.
\]

\emph{Type selectors.} Let $\alpha=(a_1,\ldots,a_l)\vdash n$ and $M_\alpha=\prod_jm_j(\alpha)!$. Consider $\Phi_\alpha(\pi)=\sum\prod_{i=1}^l\binom{|B_i|}{a_i}$ over injective assignments $i\mapsto B_i$ of distinct blocks of $\pi$. A nonzero term needs $a_i\le|B_i|$; since $\sum_ia_i=n=\sum_{B\in\pi}|B|$, every block is used and $|B_i|=a_i$. Hence $\Phi_\alpha(\pi)=M_\alpha$ if $\pi$ has type $\alpha$ and $0$ otherwise. By M\"obius inversion on the set partitions $\eta$ of the index set $[l]$, the injective sum equals $\sum_\eta\prod_{I\in\eta}(-1)^{|I|-1}(|I|-1)!\,\Sigma_I$, where $\Sigma_I=\sum_{B\in\pi}\prod_{i\in I}\binom{|B|}{a_i}$ is the sum over assignments sending all of $I$ to one block. Classifying tuples of subsets of a $b$-set of sizes $a_i$ ($i\in I$) by the size $s$ of their union gives $\prod_{i\in I}\binom{b}{a_i}=\sum_sN_I(s)\binom bs$, where $N_I(s)\in\Z$ counts such tuples with union a fixed $s$-set and $1\le s\le\sum_{i\in I}a_i$. Therefore $\Sigma_I=\sum_sN_I(s)B_s(X_\pi)$ is the value on $X_\pi$ of a symmetric polynomial of degree at most $\sum_{i\in I}a_i-1$, and the term of $\eta$ has degree at most $n-|\eta|\le n-1$. If $\alpha\ne(1^n)$, every $m_j(\alpha)\le n-1$, so $M_\alpha$ is a unit; dividing by it gives a symmetric polynomial $Q_\alpha$ of degree at most $n-1$ with $Q_\alpha(X_\pi)=1$ for $\type(\pi)=\alpha$ and $0$ otherwise. Put $Q_{(1^n)}=1-\sum_{\alpha\ne(1^n)}Q_\alpha$.

\emph{Prefactor.} Let $P=\sum_{k=0}^{n-1}e_k$, of degree $n-1$. Then $P(J^\pi)=\prod_i(1+J_i^\pi)-e_n(J^\pi)=H_\pi$, because $e_n(J^\pi)=J^\pi_1\cdots J^\pi_n=0$. Hence $f_\alpha=PQ_\alpha$ has degree at most $2n-2$ and $f_\alpha(J^\pi)=H_\pi Q_\alpha(J^\pi)=Q_\alpha(X_\pi)H_\pi$, so \eqref{eq:mob} gives $T_n(f_\alpha)=\mu_\alpha G_\alpha$, and Lemma~\ref{lem:G} finishes the construction.

\emph{Lower bound.} A homogeneous $f$ of degree $d$ expands into words of length $d$. Let $w=(t_1,\ldots,t_d)$ be a transitive word with product $1$. At each step the partial product lies in the group generated by the earlier transpositions, so its cycles lie inside the current components of the graph. A transposition joining two components therefore joins two cycles, lowering the number of cycles by one; transitivity needs at least $n-1$ such steps. The number of cycles starts and ends at $n$ and changes by $\pm1$ at each step, so there are as many cuts as joins, and $d\ge2(n-1)$. Hence words of length less than $2n-2$ contribute nothing to the identity coefficient. Since $R\ne0$, the identity element (the class sum $K_{(1^n)}$) has no preimage of degree less than $2n-2$.
\end{proof}

\section{Ordinary and transitive evaluations together}\label{sec:pair}
This section determines which pairs $(E_n(f),T_n(f))$ occur. Only one linear constraint links them.

\begin{theorem}[Simultaneous image]\label{th:pair}
If $(n-1)!$ is a unit in $R$, the image of $f\mapsto(E_n(f),T_n(f))$ consists exactly of the pairs $(A,B)$ of central elements with $[\sigma]A=[\sigma]B$ for an $n$-cycle $\sigma$.
\end{theorem}
\begin{proof}
A word whose product is an $n$-cycle generates a group containing that $n$-cycle, so it is transitive; hence $[\sigma]E_n(f)=[\sigma]T_n(f)$ for every $n$-cycle $\sigma$. Conversely, let $(A,B)$ satisfy the condition. Choose $f$ with $E_n(f)=A$ (Section~\ref{sec:setup}). For $\alpha\ne(n)$, the function $F_\alpha$ of Theorem~\ref{th:centre} satisfies $E_n(F_\alpha)=F_\alpha(J^{\hat1})=0$, because $\hat1$ has type $(n)$ and the marker monomial of $F(J^{\hat1};y)$ is that of type $(n)$. The elements $T_n(F_\alpha)=\mu_\alpha G_\alpha$, $\alpha\ne(n)$, have zero coefficient at the $n$-cycles and, by Lemma~\ref{lem:G}, form a basis of the central elements with that property. Since $B-T_n(f)$ has zero $n$-cycle coefficient, adding a suitable combination of the $F_\alpha$ to $f$ realises $B$ without changing $E_n(f)$.
\end{proof}

\section{Algebra generation and the integral image}\label{sec:integral}
This section shows that multiplication does not remove the obstruction of Theorem~\ref{th:centre}, and then asks what the transitive images span over the integers. The basic tool is a conjugation-orbit bound.

\begin{proposition}[Classwise divisibility]\label{prop:orbit}
Let $x$ be a product of one or more elements of $T_n(\Lambda_\Z)$. For every $\lambda\vdash n$, the coefficient of $K_\lambda$ in $x$ is divisible by $z_\lambda/h_\lambda$. Consequently every element of the unital $\Z$-algebra generated by $T_n(\Lambda_\Z)$ has coefficient divisible by $z_\lambda/h_\lambda$ at every class $\lambda\ne(1^n)$, and every element of $T_n(\Lambda_\Z)$ also has identity coefficient divisible by $z_{(1^n)}/h_{(1^n)}=(n-1)!$.
\end{proposition}
\begin{proof}
By linearity we may take $x=T_n(e_{\nu^{(1)}})\cdots T_n(e_{\nu^{(s)}})$. Fix $\sigma$ of type $\lambda$. By \eqref{eq:count}, $[\sigma]x$ counts sequences of $s$ tuples of permutations, with prescribed cycle numbers, such that each tuple generates a transitive group and the product of everything is $\sigma$. The centraliser $C(\sigma)$, of order $z_\lambda$, acts on these sequences by simultaneous conjugation. An element $c$ of a stabiliser commutes with the first transitive group, so, as in the proof of Theorem~\ref{th:centre}, it acts freely on $[n]$. It also commutes with $\sigma$, so it preserves each nonempty set $U_i$ of letters lying in $i$-cycles of $\sigma$ and acts freely on it. Hence the stabiliser order divides $|U_i|=i\,m_i(\lambda)$ for each $i$, and so divides $h_\lambda$. Every orbit has size $z_\lambda/|\mathrm{stab}|$, a multiple of $z_\lambda/h_\lambda$, and $[\sigma]x$ is a sum of orbit sizes. The scalars contribute only to the identity class.
\end{proof}

Thus $z_\lambda/h_\lambda$ divides the size of every orbit of the centraliser of $\sigma$ on transitive factorisations of $\sigma$. It plays the role of the automorphism factor familiar from Hurwitz theory.

\begin{theorem}[Unital generation]\label{th:unital}
The unital $R$-subalgebra of $Z(R\Sn_n)$ generated by $T_n(\Lambda_R)$ is the whole centre if and only if $(n-1)!$ is a unit in $R$. More precisely, for each prime $p\le n-1$, every element of that subalgebra has coefficient in $pR$ at each permutation consisting of one $p$-cycle and $n-p$ fixed points.
\end{theorem}
\begin{proof}
The subalgebra is the $R$-span of $1$ and of products of elements of $T_n(\Lambda_\Z)$. For $\lambda=(p,1^{n-p})$ with $p\le n-1$ prime, $z_\lambda=p\,(n-p)!$ and $h_\lambda=\gcd(p,n-p)$, which is $1$ or $p$. If $h_\lambda=1$ then $p$ divides $z_\lambda/h_\lambda=p(n-p)!$; if $h_\lambda=p$ then $p$ divides $n-p\ge p$, so $p$ divides $z_\lambda/h_\lambda=(n-p)!$. By Proposition~\ref{prop:orbit} the coefficient at this class lies in $pR$. If $(n-1)!$ is not a unit, some prime $p\le n-1$ is not a unit, so $pR\ne R$ and $K_{(p,1^{n-p})}$ is not in the subalgebra. If $(n-1)!$ is a unit, Theorem~\ref{th:centre} already gives the whole centre.
\end{proof}

\paragraph{The integral image.} Since $T_n(\Lambda_R)$ is the $R$-span of $T_n(\Lambda_\Z)$, the lattice $T_n(\Lambda_\Z)\subseteq Z(\Z\Sn_n)$ determines the linear image over every coefficient ring, and the unital algebra it generates determines the generated algebra. Proposition~\ref{prop:orbit} shows that $T_n(\Lambda_\Z)$ lies in the lattice
\[
 L_n=\Bigl\{\sum_\lambda c_\lambda K_\lambda:\ c_\lambda\in(z_\lambda/h_\lambda)\Z\text{ for every }\lambda\vdash n\Bigr\}.
\]
Theorem~\ref{th:centre}(ii) says that the identity coordinate of $L_n$ is attained. It is natural to ask whether $T_n(\Lambda_\Z)=L_n$. We computed $T_n(\Lambda_\Z)$ exactly for $n\le7$.

\begin{theorem}[Integral image for $n\le7$; computer-assisted]\label{th:integral}
\begin{enumerate}[label=(\roman*)]
\item For $n\in\{1,2,3,4,5,7\}$, $T_n(\Lambda_\Z)=L_n$.
\item For $n=6$, $T_6(\Lambda_\Z)$ has index $2$ in $L_6$. It consists of the elements of $L_6$ satisfying
\[
 \frac{c_{(4,1,1)}}{4}+\frac{c_{(2,2,2)}}{8}+\frac{c_{(2,1,1,1,1)}}{24}\equiv0\pmod 2 .
\]
In particular $24\,K_{(2,1,1,1,1)}\in L_6$ is not a transitive image of an integral symmetric function.
\item For every $n\le7$, $\Z\cdot1+T_n(\Lambda_\Z)$ is closed under multiplication, so it is the unital $\Z$-algebra generated by the transitive images. Its index in $Z(\Z\Sn_n)$ is $1$, $1$, $2$, $12$, $13824$, $26542080$ and $285315214344192000$ for $n=1,\ldots,7$.
\end{enumerate}
\end{theorem}
\begin{proof}[Method of computation]
Fix $\sigma$ of type $\lambda$. Only the set partitions $\pi$ coarser than the cycle partition of $\sigma$ contribute to $[\sigma]$ of \eqref{eq:mob}; each is a grouping of the cycles of $\sigma$ into blocks. The element $f(J^\pi)$ acts on the irreducible representation $\bigotimes_B\rho_B$ of $\Sn_\pi$ by the scalar $f(c)$, where $c$ is the union over the blocks of the multisets of contents of the diagrams $\rho_B$ (on each block, symmetric functions of Jucys--Murphy elements act through contents). Hence the central character formula gives $[\sigma]T_n(f)=\sum_cw_{\lambda,c}\,f(c)$, a finite sum over multisets $c$ of $n$ integers with explicit rational weights $w_{\lambda,c}$. The evaluation map $f\mapsto(f(c))_c$ is a ring homomorphism from $\Lambda_\Z$ to $\Z^{M}$, where $M$ is the finite set of multisets that occur. Its image is the span of products of the images of $e_1,\ldots,e_{n-1}$, because $e_k$ vanishes on every $c$ for $k\ge n$ (every content multiset contains $0$). The span of products of at most $m$ generators grows with $m$; once two consecutive spans agree, all later ones do, and the span is the whole image. The program computes these spans by Hermite normal forms until they stabilise, applies the weights, and obtains a $\Z$-basis of $T_n(\Lambda_\Z)$ with no bound on degree. It then compares this basis with $L_n$, verifies the indices and the congruence in (i)--(ii), and tests products of basis elements of $\Z\cdot1+T_n(\Lambda_\Z)$ for membership. All arithmetic is exact. The engine is checked against direct word expansion for $n\le4$ (Section~\ref{sec:checks}).
\end{proof}

By Theorem~\ref{th:integral}(i), for $n\le5$ and $n=7$ the image over any ring $R$ is $\bigoplus_\lambda(z_\lambda/h_\lambda)R\,K_\lambda$, and the unital algebra it generates is $R\cdot1$ plus that module. Part (ii) shows that the classwise bound of Proposition~\ref{prop:orbit} is not the whole story. All three classes in the congruence consist of odd permutations, but we have no conceptual explanation of it.

\begin{question}
Describe $T_n(\Lambda_\Z)$ for all $n$. Is its index in $L_n$ always a power of $2$, and is it $1$ for odd $n$? Is $\Z\cdot1+T_n(\Lambda_\Z)$ always closed under multiplication?
\end{question}

\section{Computational checks}\label{sec:checks}
This section lists the computer checks that accompany the proofs. They are exact, use only the Python standard library, and are run by one script, \texttt{code/verify\_transitive.py}, which writes \texttt{results/verification.json}.
\begin{enumerate}[label=(\alph*)]
\item For $n\le4$, the transitive part of the marker polynomial $F(x;y)$ of Theorem~\ref{th:centre} is computed by expanding words directly and tracking their orbits, and agrees with $\sum_\pi\mu(\pi)H_\pi\prod_jy_j^{a_j(\pi)}$ (with denominators cleared).
\item For $n\le4$, Lemma~\ref{lem:mob} holds for all $158$ ordered monomials of length at most $4$, and reversing the order of a monomial does not change its transitive image.
\item For $n\le9$, the type selectors of Theorem~\ref{th:degree} take the stated values on all $1818$ pairs of partition types, and every inclusion--exclusion term has degree at most $n-1$.
\item For $n\le4$, the class-level engine of Theorem~\ref{th:integral} agrees with direct word expansion for every $e_\nu$ of degree at most $2n-2$; the $n=3$ matrix of Section~\ref{sec:linear} is reproduced.
\item For $2\le n\le7$, no $e_\nu$ of degree less than $2n-2$ has a transitive image with nonzero identity coefficient, and those of degree at most $2n-2$ span the centre over $\Q$ (Theorem~\ref{th:degree}). For $2\le n\le6$, the pair map of Theorem~\ref{th:pair} has rank $2p(n)-1$ over $\Q$, where $p(n)$ is the number of partitions of $n$.
\item For $n\le7$, the integral image, its containment in $L_n$, the index and congruence of Theorem~\ref{th:integral}, and the closure of $\Z\cdot1+T_n(\Lambda_\Z)$ under multiplication.
\end{enumerate}
The checks support the proofs; the proofs of Sections~\ref{sec:setup}--\ref{sec:integral} do not depend on them, except Theorem~\ref{th:integral}, which is a computation.

\section{Remarks}\label{sec:remarks}
This section notes a wider use of the method and records how the note was prepared.

The M\"obius formula \eqref{eq:mob} is not specific to symmetric functions. It applies coefficientwise to any generating series of words built from commuting families of Jucys--Murphy elements on blocks. One example is the family of monotone words in which each larger letter is marked when it is first used. These words lead to the deformed monotone Hurwitz polynomials of Coulter, Do and Moskovsky \cite{cdm}. We do not pursue that application here.

\paragraph{AI-use statement.} This note was prepared with substantial assistance from AI systems, which drafted proofs, wrote the verification code and carried out internal review. The arguments were checked step by step, and every computational statement is produced by the accompanying script. The note has not been refereed.

\paragraph{Data availability.} The verification script and its output are distributed with this note.

\begin{thebibliography}{9}
\bibitem{clr} J.~Campion Loth and A.~Rattan. Centrality of star and monotone factorisations. \emph{Bulletin of the London Mathematical Society} 57 (2025), 3567--3585. \url{https://doi.org/10.1112/blms.70170}. Preprint \url{https://arxiv.org/abs/2403.08354}.
\bibitem{cdm} X.~Coulter, N.~Do and E.~Moskovsky. Integration on complex Grassmannians, deformed monotone Hurwitz numbers, and interlacing phenomena. \emph{Canadian Journal of Mathematics} (2025), 1--40. \url{https://doi.org/10.4153/S0008414X25101168}.
\bibitem{feray} V.~F\'eray. Partial Jucys--Murphy elements and star factorizations. \emph{European Journal of Combinatorics} 33 (2012), 189--198. \url{https://doi.org/10.1016/j.ejc.2011.09.035}. Preprint \url{https://arxiv.org/abs/0904.4854}.
\bibitem{gj} I.~P.~Goulden and D.~M.~Jackson. Transitive powers of Young--Jucys--Murphy elements are central. \emph{Journal of Algebra} 321 (2009), 1826--1835. \url{https://doi.org/10.1016/j.jalgebra.2009.01.004}. Preprint \url{https://arxiv.org/abs/0704.1100}.
\bibitem{jucys} A.-A.~A.~Jucys. Symmetric polynomials and the center of the symmetric group ring. \emph{Reports on Mathematical Physics} 5 (1974), 107--112. \url{https://doi.org/10.1016/0034-4877(74)90019-6}.
\bibitem{murphy} G.~E.~Murphy. The idempotents of the symmetric group and Nakayama's conjecture. \emph{Journal of Algebra} 81 (1983), 258--265. \url{https://doi.org/10.1016/0021-8693(83)90219-3}.
\bibitem{murray} J.~Murray. Generators for the centre of the group algebra of a symmetric group. \emph{Journal of Algebra} 271 (2004), 725--748. \url{https://doi.org/10.1016/j.jalgebra.2003.02.007}.
\end{thebibliography}
\end{document}
